% Faithful transcription --- original Latin, transcribed from the ORIGINAL publication in the
% Acta Eruditorum, Leipzig, June 1694, pp. 276--280: Jacob Bernoulli's "Solutio Problematis
% Leibnitiani de Curva Accessus & Recessus aequabilis a puncto dato, mediante rectificatione Curvae
% Elasticae". Scan: Internet Archive, Acta Eruditorum 1694 (id s1id13206610), leaves 0290--0294.
% \origpage uses the PRINTED Acta page numbers.
%
% NOTE ON SOURCE: this work was first drafted from the 1744 collected Opera (No. LIX, pp. 601--607)
% and has now been re-collated against, and conformed to, the 1694 Acta original. Consequences:
%  * In the Acta this is NOT a standalone "No. LIX": it is the closing section of Bernoulli's big
%    June 1694 elastica article, introduced by the running heading "Jac. B. Solutio Problematis
%    Leibnitiani de Curva Accessus & Recessus aequabilis ...". The "No. LIX" number and the Opera
%    page numbers 601--607 were the collected edition's; both are dropped here.
%  * The lettered footnotes (a)--(m) were Gabriel Cramer's 1744 editorial apparatus and are absent
%    from the Acta (removed earlier).
%  * The central differential equation is CONSISTENT in the Acta: it reads
%    $(x\,dx+y\,dy)\sqrt{y}=(x\,dy-y\,dx)\sqrt{a}$ both at its first statement and in observation 8.
%    The variant form printed in the Opera's observation 8 was an Opera misprint and is not kept.
%
% Transcription notes for the reviewer:
%  * Author notation is kept faithfully (HOUSESTYLE R3): doubled letters for squares
%    (aa, zz, xx, yy, tt, uu), the ":" ratio/division sign, the classical proportion sign "::"
%    ("as ... so ..."), "&c." for et cetera, the "&" ampersand, archaic spelling. Typography is
%    normalized: long-s, ligatures, small caps, and the superscript-"o" of "No." are set as plain
%    letters; the author's parenthetical insertions and value-substitutions use round parentheses.
%  * The five compounded ratios (p. 278) are set as the Acta prints them: "Elem. periph. Be" (not
%    "Ge"), with NO "(el)"/"(ek)" substitutions (those were Cramer's cross-references to his deleted
%    note c), and the proportion sign "::" for "as ... so ...". The proportion sign "::" is used
%    throughout --- both the five ratios and the running proportion chain ("unde facta divisione
%    ...") --- while single-colon ":" stays for a ratio/division and "=" is reserved for a genuine
%    equality (e.g. "Ergo alphabeta = Aomega"), matching the Acta literally.
%  * Reviewer flags: (i) the little isochronal arc is set "A\omega"; the Acta glyph is a curly
%    varpi-like form, read here as omega (the defined point) --- confirm. (ii) obs. 7 prints
%    "horizantali" (an a-spelling; obs. 3 has "horizontalem"): kept as printed (R4).
\documentclass{article}
\usepackage{readmasters}
\begin{document}

\origpage{276}
\section*{Jac.\ B. Solutio Problematis Leibnitiani}
\subsection*{de Curva Accessus \& Recessus aequabilis a puncto dato, mediante rectificatione Curvae Elasticae}

\origpage{277}
Felicitatem inventi praecedentis commendare potest solutio elegantissimi Problematis
\emph{Leibnitiani}, de invenienda \emph{Linea, per quam descendens grave aequaliter aequalibus
temporibus a dato puncto recedat, vel ad illud accedat}: quod laudatissimo suo Auctori ita placuit,
ut non tantum ad ejus tentamen Amicos pariter \& Adversarios aliquoties in \emph{Actis} provocarit
(vid.\ \emph{Anno}~1689, \emph{Apr.\ p.}~198, \& 1690. \emph{Maj.\ p.}~229, \& \emph{Jul.\ p.}~360,)
sed \& ipse quoque strenue in illo desudarit, testibus nonnullis Curvae proprietatibus, quas
\emph{Gallorum Diario} inseri curavit, ut ex relatione \emph{Fratris} habeo, qui tamen illas nominare
mihi non potuit. Multus etiam in eodem fuit ipse \emph{Frater}, dum \emph{Parisiis} ageret, sed omni
sua methodo Tangentium inversa aliud nil effecit, quam ut Problema ad hanc aequationem differentialem
reduceret: $(x\,dx + y\,dy)\sqrt{y} = (x\,dy - y\,dx)\sqrt{a}$, ad quam tamen etiam nullo labore
pervenitur. Plenariam vero illius solutionem, sive constructionem, nec ipse nec quisquam alius dare
potuit. Nuper cum praecedens Schediasma pararem \emph{Lipsiam}, partemque inventi cum figuris ipsis
jamjam chartae consignassem, scrutinium hujus Lineae suscepi occasione novorum Theorematum, quae ibi
dedi pro Radiis circulorum osculantium in Spiralibus (quorum tamen investigatio \& ipsa incidenter
tantum fiebat; praecipua enim inventi capita sub calamo demum mihi nascebantur) mox etiam post pauca
conamina optatum finem consecutus sum; ubi hoc perquam singulare mihi accidit, ut absoluta analysi
animadverterem, Curvam Elasticam, quae qualemcunque tantum occasionem huic tentamini dederat, illam
ipsam esse, cujus rectificatione altera posset construi. Nam quanquam idem liceat exequi mediante
quadratura spatii alicujus Algebraici, alterum tamen construendi modum praeferendum censeo, tum quod
generaliter facilius in praxi rectificentur curvae, quam quadrentur spatia, tum praesertim quod ipsa
natura (si quidem convenientem tensionis legem observet) illum praescripsisse videatur.

\emph{Constr.} Descripta itaque Curva Elastica tertiae constructionis (\emph{Fig.}~3) caeterisque
positis ut prius, applicetur in utroque quadrante circuli $BL$, $LG$, recta $\varepsilon\zeta$
aequalis ipsi $Ag$ seu tertiae proportionali ad $AB$ \& $AE$ (hic in dextro tantum quadrante ducta
apparet, ut linearum confusio vitetur) \& ex ducta $A\varepsilon$ (productaque si sit opus)
abscindatur $A\alpha$ aequalis tertiae proportionali ad rectam $AB$ \& portio-

\origpage{278}
nem Curvae $AQ$ in sinistro, vel $\varphi RQ$ in dextro quadrante; eritque punctum $\alpha$ in Curva
quadam $A\chi\omega\eta$, ita comparata, ut grave per illam latum (si prius ex altitudine $iA$
deciderit) aequalibus temporibus aequaliter ad punctum $A$ accedat, aut ab illo recedat. Cujus
constructionis demonstrationem (ut Celeb.\ Problematis Auctorem, aliis occupatissimum, examinandi
labore sublevem) hic appono.

\rmfigure{figures/fig-3.png}{Fig.~3.}{Fig. 3: the elastic curve, the isochronal curve
$A\chi\omega\eta$, and the auxiliary circle, quadrants and construction lines used in the
demonstration.}

\emph{Dem.} Descriptus intelligatur centro $A$ arculus $\beta\delta$, abscindens ex applicata
$A\alpha$ Elementum $\alpha\beta$, \& ex Curva Elem.\ $\alpha\delta$, cui concipiatur in vertice
portiuncula isochronos $A\omega$. Jam Ratio $\beta\delta$ ad $\alpha\beta$ composita est ex
sequentibus 5 rationibus:

\[
\begin{aligned}
\beta\delta : \text{Elem. periph. } B\varepsilon &:: A\alpha : A\varepsilon. \\
\text{Elem. periph. } B\varepsilon : \text{Elem. } \varepsilon\zeta &:: A\varepsilon : A\zeta, \qquad \text{ex natura Circuli.} \\
\text{Elem. } \varepsilon\zeta : \text{Elem. } AE &:: 2AE : AB, \qquad \text{per constr. \& princ. calcul. in fin. parvorum.} \\
\text{Elem. } AE : \text{Elem. } AQ\ (\varphi RQ) &:: A\zeta : AB, \qquad \text{ex nat. curvae } AQ, \text{ ut collig. ex coroll. 2. constructionis tertiae praeced.} \\
\text{Elem. } AQ\ (\varphi RQ) : \alpha\beta &:: AB : 2AQ\ (2\varphi RQ), \qquad \text{per constr. \& princ. calc. inf. parv.}
\end{aligned}
\]

unde facta divisione per communes rationum antecedentes \& consequentes,
$\beta\delta : \alpha\beta :: A\alpha : AB + AE : AQ\ (\varphi RQ,)$ \&
$\beta\delta^{2} : \alpha\beta^{2} :: A\alpha^{2} : AB^{2} + AE^{2} : AQ^{2}$ vel $\varphi RQ^{2}$
($\varepsilon\zeta : A\alpha$, per constr.)
$:: A\alpha^{2} \times \varepsilon\zeta : AB^{2} \times A\alpha ::
A\alpha \times \varepsilon\zeta : AB^{2} ::{}$ (ob similitudinem Triangulorum $A\alpha\gamma$,
$A\varepsilon\zeta$) $A\varepsilon \times \alpha\gamma : AB^{2} :: \alpha\gamma : AB\ (Ai)$,
componendoque $\alpha\delta^{2} : \alpha\beta^{2} :: \alpha\gamma + Ai : Ai$; \& proinde
$\alpha\delta : \alpha\beta :: \sqrt{\alpha\gamma + Ai} : \sqrt{Ai} ::{}$ celeritas acquisita in
$\alpha$ : celer.\ acquis.\ in $A ::{}$ (ex natura descensus gravium) $:: \alpha\delta : A\omega$ (per
hyp.\ quia aequalibus tempusculis transiri supponuntur). Ergo $\alpha\beta = A\omega$. Ergo grave
per hanc Curvam latum aequaliter aequalibus temporibus a puncto $A$ recedit. Q.~E.~D.

Quae hic praecipue observanda veniunt, sunt sequentia:

\textbf{1.}~Si grave altius vel humilius, quam ex $i$, decidere supponatur, puta ex $\tau$, non opus
est nova Elastica: Sufficit haec una omnibus Curvis isochronis describendis. Omnes enim inter se
similes sunt, determinanturque, faciendo tantum, ut $Ai$ ad $A\tau$, sic $A\alpha$ ad aliam
$A\alpha$ abscindendam ex eadem $A\varepsilon$. Estque ipsa curva Schematis nostri constructa pro
gravi tanquam ex $\tau$, non $i$, decidente.

\textbf{2.}~Sed \& infinitae dantur aliae, quas grave ex eadem altitudine

\origpage{279}
$\tau A$ delapsum ita permeare potest, ut aequabiliter ad punctum aliquod accedat vel recedat; at
illud continuo diversum ab $A$; nec nisi per unicam curvam ex eadem altitudine, respectu ejusdem
puncti, accessus vel recessus aequabilis effici potest.

\textbf{3.}~Curva nostra Isochronos initium capit in ipso puncto $A$, cujus respectu motus aequabilis
est, \& terminatur in puncto $\eta$ ejusdem cum illo altitudinis; utraque vero extremitate tangit
horizontalem rectam $A\eta$. Unde flexum contrarium habere constat. Quod si ex altera parte
statuatur huic similis \& aequalis curva $A\lambda\theta\psi$, \& grave dimittatur ex $\eta$ cum
celeritate, quam acquirere potest cadendo ex altitudine aequali ipsi $\tau A$, illud primo
descendendo ad $\omega$, mox reascendendo, aequabiliter ad $A$ accedet, indeque pergendo continue per
alteram $A\lambda\theta\psi$ ab eodem puncto $A$ eadem ratione elongabitur.

\textbf{4.}~Tangens Curvae in locis intermediis, velut $\alpha$, reperitur, si ducta $A\mu$
perpendiculari ad $A\alpha$, fiat, ut $\sqrt{A\tau}$ ad $\sqrt{\alpha\gamma}$; ita $A\alpha$ ad
$A\mu$; ducta enim $\alpha\mu$ curvam tanget: quod vel ex usu, quem praestare debet curva, facile
infertur.

\textbf{5.}~Recta $A\eta$ quadrupla est rectae $A\theta$.

\textbf{6.}~Si in Elastica punctum $Q$ assumatur tale, ut tertia proportionalis ad subtangentem $Pp$
\& rectam $AB$ aequetur curvae $AQ$, invenietur ejus ope punctum $\chi$ portionis $A\chi\theta$, quod
omnium maxime distat a perpendiculari $A\theta$.

\textbf{7.}~Sin punctum $Q$ ponatur tale, ut tangens $Qp$ aequetur curvae $\varphi RQ$, illo mediante
habebitur in altera punctum $\omega$, omnium infimum \& remotissimum ab horizantali $A\eta$.

\textbf{8.}~Tandem \& hoc advertendum est, quod ex propositi Problematis constructione proclive jam
colligere sit, quales suppositiones faciendae, ad reducendam aequationem
$(x\,dx + y\,dy)\sqrt{y} = (x\,dy - y\,dx)\sqrt{a}$. Nam si loco $x$ \& $y$ assumantur duae aliae
indeterminatae $t$ \& $z$, ponendo $ay = tz$, \& $ax = t\sqrt{aa - zz}$; prodibit aequatio, in qua
separari possunt indeterminatae $t$ \& $z$ cum suis differentialibus a se invicem; quod sufficit ad
constructionem deinde, saltem per quadraturas, expediendam. Duo enim in hoc calculo adhuc praecipue
desiderari videntur, quae si semper fieri possent, omnia reperta essent; unum ut differentialia
secundi altiorumve generum ad differentialia primi reducantur; alterum, ut in aequationibus
differentialibus primi

\origpage{280}
generis indeterminatae, si invicem permixtae sint, a se mutuo separentur, ut unaquaeque cum sua
differentiali peculiarem aequationis partem constituat, h.\ e.\ ut reperiatur methodus inversa
tangentium. Pro utroque dedi Regulas quasdam (pluresque sine fine dare possem) similes illis, quas
\emph{Frater} \emph{Parisiis} apud \emph{Marchionem} Hospitalium deposuit, quibusque methodum inesse
initio existimarat. At statim sensi, illas non continere nisi artificia quaedam particularia, quae
methodum appellare non ausim, utpote qualem non magis dari posse arbitror, quam dari potest
universalis methodus pro construendis Aequationibus Algebraicis quorumvis promiscue graduum. Illa
enim (cujus olim in \emph{Actis} mentio facta est) qua pono indefinite $0 = a + bx + cy + exy$, \&c.\
\& quaero deinde ejus ope lineae tangentem, vix aliter quam in speculatione succedere videtur. Quod
dictis fidem faciat, hoc esto, quod praesens Problema, quanquam ut apparet haud admodum difficile,
nullius regulae methodive hactenus repertae ope solvi a quoquam potuerit.

\end{document}
